% 第七章 雍正争位与禁教
% 原书页码：249；标题
\section{雍正争位与禁教}
\par
% 原书页码：249；段落
雍正五年四月初八（1727年5月24日），雍正收到了广东发来的%
消息，新来的葡萄牙使团已经离开广东赶来北京了。这一天正好是释%
迦牟尼佛的生日，雍正在朝会上公开评论了天主教和佛教：%
\par
% 原书页码：249；段落
{\bookkai 今日为佛诞之期，恰遇西洋国使臣上表称贺，两事适然相值，尔}%
{\bookkai 等在廷诸臣奏事之暇，偶将朕意宣谕尔等知之。……朕意西洋立教之}%
{\bookkai 初，其人为本国所敬信或者尊之如天。倘谓立教之人，居然自称为天}%
{\bookkai 主，此理之所无者也。释氏原以清净无为为本，以明心见性为功，所}%
{\bookkai 以自修自全之道，莫善于此。\NoteRef{7}{1}}%
\par
% 原书页码：249–250；段落（跨页续段）
雍正在朝会上公开对着群臣臧否宗教，是很少见的，特别是他评点的%
内容和时政又没有什么关系。雍正大概也知道这有点不合常规，所以%
还专门解释了一下，他只是在大臣们“奏事之暇”，抽时间说几句。%
很明显，雍正这段话说明他对天主教中像“天主”这样最基本的概念%
都是不认同的。尽管他措词委婉，“此理之所无者也”，但明白人都能%
听出，雍正是从根本上否定了天主教。同时，雍正也讲出了他心中最%
好的宗教，那就是释氏的佛教。在朝堂上，面对朝臣，这样公开的一%
贬一褒，雍正的态度不能再明确了。%
\par
% 原书页码：250；段落
其实，早在三年以前，雍正元年都还没过完，他就废除了康熙的%
容教诏令，禁止传教士在中国地方上继续传教。传教士只能留在京城%
为朝廷做天文观测这样的技术性工作。当然，北京的传教士也没闲着。%
从雍正禁教起，他们就在北京找各种关系，希望说服雍正取消禁令。%
雍正也明白这个在路上的葡萄牙使团，大概又是来求情的。在葡萄牙%
使团来华前两年，雍正已经拒绝过教皇派来的使团。\NoteRef{7}{2}所以雍正在收到%
葡萄牙公文以后，就在朝堂上公开评价了天主教，并让官员们把他的%
意思带给葡萄牙使团。%
\par
% 原书页码：250–251；段落（跨页续段）
北京传教士烦恼的不仅仅是他们的宗教这样被皇帝公开羞辱，%
最让他们沮丧的是他们看不到一点儿希望。康熙去世后，他们在北%
京的状况江河日下，一步步被排挤出了北京的权贵圈。雍正登基后，%
第一件事就是把他自己府上的和信任的人，换去管理内务府。传教%
士作为康熙内务府的一部分，与康熙其他近身奴才遭遇基本一样，%
都被逐出了宫，不再属于内务府。据马国贤神父回忆录所记，雍正%
元年管理内务府的是康熙的十七皇子，他把神父们召集到一起，通%
知他们以后就不再归内务府管，也不算内务府的人了。不归内务府%
管也就是被剥夺了内务府身份，最直接的一个结果就是以后他们得%
不到皇帝的专门召见，不能人宫了。皇十七子还专门强调了这一点。\NoteRef{7}{3}%
能入宫面圣是了不起的特权。这个传教士享有了超过半个世纪的特%
权，就这样被取消了。最让传教士们失望的是，雍正皇帝连表现机%
会都没有给过他们，就直接把他们赶出了内务府。没有了内务府的%
身份，他们就只是一群懂天象观测的术士，地位一下子就跟市场上%
掐八字算卦的没多大区别了。%
\par
% 原书页码：251；段落
从被赶出内务府，到雍正全面禁止天主教，中间大概相隔了九个%
月时间。这期间，传教士已从朝野上下对他们态度的转变中，预感到%
了禁教的前景。被逐出内务府对传教士来说是釜底抽薪，一下断了他%
们在大清朝野存在的根本，朝野上下也都明白这其中的意味。传教士%
很真切地感受到了各方面势力对他们态度的变化。很快，许多权贵都%
开始回避他们，不和他们来往了。%
\par
% 原书页码：251；段落
雍正继位的时候已经四十五岁，思想、性格都已经完全成熟。对%
于传教士来说，他们没有任何机会影响、改变雍正。这一点，跟当年%
七岁继位的康熙不一样。传教士那时接触到的康熙，正处在学习阶段。%
他们成功地用他们的知识来引导康熙，让康熙对西洋知识、器物都发%
生了兴趣。可惜的是传教士的一切在雍正那里都没有吸引力。雍正%
二十岁开始就喜欢和喇嘛、和尚、道士往来。在雍正还是皇子的时候，%
传教士也没有特别花过功夫去接近雍正，双方尽管没有矛盾，但也没%
有特别的关系。前章讲到，传教士在康熙朝受到的厚待得益于他们和%
康熙之间的关系，那么雍正朝他们受到的冷遇，也是他们和雍正关系%
的体现。传教士和康熙各个皇子建立关系的黄金时段是康熙晚年废掉%
太子以后，也就是康熙的最后十年。在这段时间里，为什么传教士会%
错过雍正，没有和他建立起良好的关系呢？这就不得不说到康熙晚年%
的储位斗争。%
% 原书页码：252；标题
\subsection{康熙晚年的皇位争斗}
\par
% 原书页码：252–253；段落（跨页续段）
康熙晚年的储位之争大概是他一生中最感棘手的问题。涉及的人%
都是自己的儿子和亲人，手心手背都是肉，康熙很难恰到好处地处理。%
满人的传统，大多是不立太子的。努尔哈赤和皇太极去世的时候，都%
没有立太子。但也正是因为没有立太子，二人死后满人内部都有剑拔%
弩张的争位戏码上演。立太子是汉人朝廷的传统，也是皇位和平有序%
传递的一个措施。1675年，康熙听了汉人臣子的意见，在他二十一岁%
的时候，就册立了当时才两岁的皇二子为太子。太子长大后，中年的%
康熙也确实体会到了立太子的好处。康熙几次塞外远征和南巡，都是%
太子坐镇北京，管理朝廷日常运作，让康熙可以放心远走。但问题是，%
康熙走后，太子在北京朝廷中体验着当皇帝的感觉，康熙回来，他又%
回到太子的角色上去。太子等待皇位的时间是漫长的，在当了三十三%
年的太子以后，还是在1708年被废了。康熙在颁布的废太子诏中，%
从太子“秉性乖戾”说起，给出了一大堆太子不好的理由。不过，最%
本质的问题应该还是康熙察觉到了太子想登位，他自己感受到了威胁。%
但是废黜太子并不是解决皇储问题的办法。康熙很快发现这一举动反%
而激化了皇储问题中隐藏的各种矛盾。太子位空缺让他众多儿子看到%
了机会，一个个蠢蠢欲动，同时整个朝廷的各派势力也通过支持不同%
的皇子在暗中角力。康熙的大皇子甚至对康熙说，如果杀废太子下不%
了手，他可以出手。康熙一直敬重的舅舅佟国维推荐立皇八子为太子，%
康熙不同意，也让舅舅不要掺和这事。佟国维则直言康熙要让自己闭%
嘴的话，除非把自己给杀了。遇到这种家务事，舅舅又这样讲话，康%
熙应该是很烦恼的。历史记载中他和舅舅佟国维的争论只是只言片语，%
但从后来康熙的反应看，二人应该不是简单的几句争吵而已。康熙幼%
年就丧父丧母，他小时候，母亲家里的两个舅舅佟国纲、佟国维一直%
作为侍卫守在他身旁。他们之间是有深切亲情的。传教士记录中，康%
熙无论朝堂上还是私下里，都是亲切地直呼二人为“舅舅”。而佟国%
纲又在跟随康熙追击噶尔丹时战死疆场，佟国维就是康熙晚年最亲的%
长辈了。按正常情况，佟国维1719年去世，康熙会赐给他这个舅舅%
很高规格的葬礼。但实际上，康熙连谥号都没有给他。最后佟国维能%
有谥号，全靠他的儿子隆科多辅佐雍正有功。雍正上台第一年为了拉%
拢隆科多，就先给佟国维补上了谥号。这说明康熙晚年和他的舅舅为%
储位问题应该产生过很大的分歧。%
\par
% 原书页码：253；段落
康熙废黜太子不久，便发现太子身上的许多问题其实是大皇子从%
中搞鬼的缘故。这让康熙对废太子的决定有些后悔。权衡各种利弊之%
后，康熙又在一年后，也就是1709年复立了太子，希望以此结束皇%
储的纷争。太子尽管恢复了，康熙对太子的猜疑还是没法去除。康熙%
本来有南巡的雅兴，1699年、1703年、1705年、1707年数次南下，%
但经过1708年、1709年废太子后又复立太子，康熙至死都再没有长%
时间远离北京南下巡游。1712年，康熙再次废太子，公开的理由还是%
跟第一次废太子时差不多，就是太子品行等各方面都不配做一个君主。%
废太子一直被软禁，直到雍正登基后的第二年去世。康熙废太子以后，%
拒绝了所有立太子的建议，这也就开启了历史中被津津乐道的九子夺%
嫡时期。%
\par
% 原书页码：254；段落
北京传教士既是康熙宫内人事起伏的见证者也是参与者。1675%
年，当两岁的太子刚刚被立以后，传教士就开始了长远计划，努力%
接近太子和他的势力圈。这个圈子的核心是索额图，是太子母亲的%
叔父。传教士一直和索额图保持着密切的私人关系。很多传教士的%
书信中都提到过索额图，称他为“索三”，因为他是索尼家的第三子。%
索额图和传教士的往来在铲除鳌拜之前就开始了。传教士能得到康%
熙颁布的容教诏令，就是索额图帮传教士打通了礼部的各个关节。%
有了索额图的这层关系，传教士在太子还未成年之时就在太子的核%
心圈子中了。从北京传教士发回欧洲的信中可以看出，他们对太子%
赞誉有加，而且认为他能延续康熙的容教政策。1706年，教皇特使%
多罗参加了康熙的辞别会后，还去了太子那边。多罗对太子的印象%
很好，认为太子会继续容留天主教。\NoteRef{7}{4}对传教士来说，1712年太子再%
次被废也是他们的重大损失，这意味着他们三十多年来在太子身上%
付出的心血付之东流了。他们需要重新猜测谁会是新的皇帝，并和%
可能的人选建立关系。%
\par
% 原书页码：254；段落
康熙一朝，传教士前后来去已历三代。他们跟所有的皇子都有交%
集。康熙学习了西方知识以后，让传教士把许多西学科目首先翻译成%
了满文，教他的儿子和皇族子弟。在康熙晚年争位的儿子中，大概可%
以分为四派：%
\par
% 原书页码：254；段落
第一派是废太子一派，以前的太子党还在寄希望废太子能有再被%
复立的一天。%
\par
% 原书页码：254–255；段落（跨页续段）
第二派是以皇三子为中心的一派。皇三子以学问大闻名，康熙把%
跟学问相关的项目都交给皇三子去办。比如著名的《古今图书集成》%
就是皇三子在主持。还有跟西学相关的天文、数学等项目也是他在负%
责。皇三子自已就是有学识的，他带着一帮满人子弟，通过他们自己%
的天文观测数据，竟然发现了传教士所讲天文学中的漏洞。在1716年%
前后，皇三子、十五子、十六子一起上书康熙，报告根据他们从1714%
年开始持续三年对夜晚的观测，发现传教士讲的天文理论有问题，有%
误差。\NoteRef{7}{5}对于皇子们发现的问题，传教士是承认的，只是说他们给的数%
据是目前最好的数据。有新来的传教士，康熙也是派皇三子去测试来%
者是否有真学问。在欧洲大学讲授数学的耶稣会传教士戴进贤来到北%
京，康熙就是派皇三子去测试的。\NoteRef{7}{6}%
\par
% 原书页码：255；段落
皇三子在以前争位历史中被提到的不多，但是皇三子在当时是一%
个不可小视的存在。研究雍正皇帝的学者黄培就提出，康熙最中意的%
皇子是皇三子，因为根据已知康熙晚年的满汉奏折来看，康熙联系最%
频繁、交办事务最多的皇子就是皇三子。\NoteRef{7}{7}后来雍正登位以后，第一年%
就把皇三子放逐去为康熙守陵，这也反映出雍正对皇三子的忌惮。康%
熙是否有心传位于皇三子，已无从考证。这里要强调的是，研读历史%
必须要站到当时人的环境中去体会。当时所有的人，包括传教士，都%
不知道康熙会传位给谁，他们只能猜测。那么他们猜测的凭据是什%
么？很重要的一点，就是看康熙把重要的任务交给哪个皇子，平时和%
哪个皇子走得近，往来的多。从这一点看，皇三子确实和康熙来往频%
繁。从现存的满文奏折上看，就特别明显。皇三子管着大大小小很多%
事，经常出现在奏折中。%
\par
% 原书页码：255–256；段落（跨页续段）
争位中第三派势力是各种史书、小说、电影、电视都经常提到的%
皇八子一系。这一派还包括皇九子、皇十子和皇十四子。在1708年%
废太子以后，皇八子确实有很高的呼声。八旗中支持皇八子的权贵众%
多，他在皇子内部也呼声很高。康熙的舅舅佟国维就支持皇八子。后%
来朝臣发现康熙铁了心不会传位皇八子，又把目光集中在这一派中的%
皇十四子。各种野史中盛传康熙有意传位给皇十四子的说法，并非空%
穴来风。康熙在晚年把最重要的西北军务交给了皇十四子。康熙一生%
最在意的就是平定西北，消灭还没有归顺的蒙古部落。皇十四子有勇%
有谋，在处理青海、西藏等问题上让康熙很满意。只是皇十四子在西%
北的历史功绩后来在历史中被雍正抹去了。%
\par
% 原书页码：256；段落
康熙晚年，皇子内部就认为康熙让皇十四子领兵去西北，就说明%
他有意把皇位给皇十四子，同时也用这个任务来最后考察皇十四子。%
对于这一点，当时人应该都知道。传教士穆经远被抓以后的审讯记录%
里面提到这事。穆经远承认他确实听到过皇子之间说康熙让皇十四子%
领兵就是中意他的表现。\NoteRef{7}{8}%
\par
% 原书页码：256；段落
同时，约翰·贝尔在他的记录中提到，他当时在北京听说，皇位%
是要传给皇十四子，同时也说皇十四子正在处理当时最麻烦的西北军%
务问题。约翰·贝尔所在的彼得大帝使团是1721年3月离开北京的，%
也就是在康熙去世前大概一年半时间。他这条记录应该算是一个跟大%
清内部没有任何利益关联的人留下的一条珍贵的历史记录。\NoteRef{7}{9}争夺皇位%
的历史情形复杂，这里只能讲到一个脉络。就本书关注的传教士来说，%
他们也跟众多朝臣一样，把宝押在了这一派上。%
\par
% 原书页码：256–257；段落（跨页续段）
第四派就是最后登位的雍正一派。雍正一派中还有皇十三子，这%
也是历史中常常提到的人物。在上面四派中，传教士可以说跟前面三%
派都有很好的关系。唯独在皇四子雍正这里，传教士基本没有下过功%
夫。为什么在废太子以后漫长的十年中他们会恰好错过雍正？要真正%
理解这个问题，必须站在传教士的角度思考。站在传教士角度最重要%
的一条，就是要明白他们是历史中的人物，无法知道历史后来的演进，%
不知道最后皇四子会登位。他们所有的决定，都是基于他们当时所观%
所感，做出他们认为的最好选择。%
% 原书页码：257；标题
\subsection{雍正登位中的谜团}
\par
% 原书页码：257；段落
康熙是1722年冬在京郊打猎时染了风寒，突然过世的。根据传教%
士记录，康熙其实在前一年身体不佳，而1722年身体已经恢复。康%
熙1722年秋已经去热河木兰打过一次猎了，冬天这次是再次出京打%
猎。这么频繁围猎也说明康熙身体没有什么问题，起码他自已没有感%
觉到有什么问题。所以康熙并没有特别准备过传位事宜，只是身体抱%
恙以后，在病床上匆忙开始了继承人选择。根据清史官方说法，康熙%
是口谕传位给皇四子的，给出的理由是皇四子“人品贵重，深肖朕躬”，%
就是说雍正品行好，康熙认为他像自己。正史中还说，传位时还有康%
熙的几位皇子和九门提督隆科多在场。%
\par
% 原书页码：257–258；段落（跨页续段）
从官方史料记载来看，传位皇四子并没有什么疑问，但雍正篡位%
的说法一直存在。这个说法到清朝灭亡，所有宫廷文档公开以后，变%
得更盛行。尽管没有铁证，但疑点很多。通过各种清官方记录比对，%
学者发现雍正时期删除了许多康熙时期的记录，而且删除的记录正好%
跟雍正的继位有关联。民国初年清史学者孟森对雍正继位的合法性做%
过深人考证，其中提到一点怀疑，这一点也是雍正朝时就被提出的问%
题：为什么雍正两个最亲的人，都对他登位不满？这两个人一个是雍%
正的母亲，一个是皇十四子。就宫廷常态来说，皇帝儿子众多，皇子%
之间一般最亲近的是同父同母的兄弟。雍正和皇十四子就是同父同母。%
在康熙驾崩之时，皇十四子正在西北处理西藏事务，过了几个月才回%
来，回来后就和雍正闹翻了，被雍正发配到城外居住。更诡异的是，%
雍正的母亲也对雍正继位没有半点高兴的样子，在雍正把皇十四子发%
配到城外以后，雍正的母亲竟然以死抗争，撞柱而死。这个“撞柱而%
死”当然在正史中不会说到，提出这个的是雍正时代的曾静。他列举%
了雍正十大罪状，游说川陕总督岳钟琪谋反，其中一条就是雍正逼死%
生母，是大不孝。至于雍正生母是否真像曾静所写是撞柱而死已不可%
知，但雍正生母对雍正登位不满，和雍正关系很差，这在很多历史记%
录中都能看到。所以为什么雍正当皇帝，连他母亲和亲兄弟都不满？%
这不得不说是一个很大的疑点。对于曾静的指控，雍正曾经公开逐条%
批判。雍正把逐条批判，以及曾静被抓后写的悔过书，一起结集出版，%
取名《大义觉迷录》。雍正这样一条条批判，相当于把当年的各种传言，%
以及宫廷矛盾记录下来，公布于众。但问题是雍正写的东西并没有起%
到他想要的效果，反而暴露出更多的疑点。因为他的记录又和当时其%
他正史记录有出人。%
\par
% 原书页码：258–259；段落（跨页续段）
乾隆初年，为了永远平息相关争议，乾隆又花大力气在全国范围%
内收回销毁了这本书。清末留日学生在日本发现了这本书，使其重见%
天日。\NoteRef{7}{10}孟森先生的考证很多就取自《大义觉迷录》，指出了很多细%
节与公布的正史之间的矛盾。\NoteRef{7}{11}后世怀疑的逻辑总结起来就是：如果%
继位没有问题，为什么会有那么多细节上的出入？而且是在雍正的御%
用学者编辑校对了正史各种记录以后，都还有那么多出人。也正因为%
存在各种疑点，雍正的继位问题成为后世每一位研究雍正历史的学者%
都绕不开的问题。当然这些疑点也为后世各种文学创作留下了想象的%
空间。%
\par
% 原书页码：259–260；段落（跨页续段）
了解了雍正继位的疑点，也才能理解传教士在雍正继位之前的%
言行。第一，传教士跟当时大多数的朝臣一样，并没有预测到雍正%
能够继位。这也是为什么传教士在雍正继位之前并没有专门花精力%
接近雍正，跟雍正私下建立起关系的原因。传教士内部其实分析过%
谁有可能继位，但他们并没有想到雍正。他们也跟当时的主流分析%
一样，把重心放在了皇八子和皇十四子一派上。当时，皇九子、皇%
十子在京城负责给在西北的皇十四子通报朝廷中的讯息。而传教士%
穆经远就属于皇九子一系，参与了康熙晚年皇子间的争位斗争。清%
亡以后，民国年间，雍正上位后处置皇八子允禩和皇九子允禟一派%
的案子就是当时的研究热点。当时故宫博物院把宫中发现跟案件相%
关的审问材料结集出版，取名为《允禩允禟案》。其中的第一份口供%
材料就是穆经远在雍正初年被抓以后录的。\NoteRef{7}{12}在罗马还存有一份当年%
刑部根据口供写的中文判决书（见附录四：穆经远判决书）。\NoteRef{7}{13}口供中%
他坦然承认和皇九子的密切关系，开头就说，皇九子对他好，是“天%
下共知”的。皇九子被雍正流放，穆经远就跟着他。单看口供，根%
本看不出他是一个西洋神父，反而处处显示出他是一个忠于主子的%
满人奴才。录口供的人还问穆经远，抓你的那天，你对着皇九子的%
房子跪地磕头是什么意思？这些如果还是小细节的话，那么穆经远%
口供里还说，他在康熙末年私自联系过年羹尧，给年羹尧送过礼物，%
这就是大事了。穆经远在康熙晚年还向康熙推荐立皇九子为嗣，当%
然这被康熙呵斥了，说他不懂规矩。这些事情都说明传教士是参与%
到了康熙晚年皇子夺位政治中的，他们心中有支持的候选人。不过%
可惜的是，他们的候选人没能登上大位。%
\par
% 原书页码：260–261；段落（跨页续段）
雍正的继位出乎朝廷上下意料，加上当时还有各种对他继位的流%
言，这些决定了雍正继位头几年的重点任务是逐步清除异己。雍正继%
位后的第一步就是更换内务府里面的人。那些常年跟在康熙身边的奴%
才都被换掉了。赵昌是康熙身边一个重要的贴身奴才。很多奏折上都%
能看到他的名字（注意赵昌是满人中的包衣奴才，不是太监，康熙还%
在其生病时，劝他禁房事）。传教士和他的私人关系也很好。雍正登基%
后，赵昌就被赶出宫，关在一个监狱里面。\NoteRef{7}{14}所以传教士们被赶出内%
务府，其实也不是雍正专门针对传教士的措施，只是雍正重新配置内%
务府人员的一个操作。后来乾隆上台以后，也把内务府里雍正养的和%
尚道士驱逐出宫，并且警告他们出去后不要乱说话。内务府就是管理%
皇帝家务事的部门，皇帝上任后，这里面的好多人物都要置换，属于%
常规操作。雍正朝由于有继位风波，外加雍正疑心大，当时被赶出内%
务府的人很多。内务府调整完后，雍正的第二步就是处理让他感到威%
胁的众兄弟们。雍正1722年冬继位，几个月后，就把皇三子发配去%
给康熙守墓，把皇九子发配到西宁，皇十子被移送到蒙古。皇十四子%
在康熙死后几个月赶回北京后，被要求在北京城外居住。当时皇八子%
在京城势力很大，人脉很广，雍正暂时没有动他，反而任用了他。雍%
正稳住大局后，在雍正四年正月把皇八子圈禁于宗人府，当年10月%
皇八子就死了。到皇八子死为止，可以说雍正前四年的执政重点都在%
打击异己、巩固皇位上。\NoteRef{7}{15}雍正和康熙以及后来的乾隆不一样，他在%
位期间没有长期离开过北京，也从来没有下过江南。雍正似乎对他的%
江山一直到死都有几分不安心。%
% 原书页码：261；标题
\subsection{传教士也尽力了}
\par
% 原书页码：261；段落
雍正元年秋（1723年10月），闽浙总督满保上了个折子，细数天%
主教在福建的问题，建议全国禁教。这个折子从内容上看，并没有多%
少新意，还是老生常谈，总体就是说天主教教义与传统观念有诸多冲%
突，不适合在中国传播等等。折子建议除有技术的北京传教士继续在%
朝廷留用以外，地方各省都应该禁止传教士出入传教。%
\par
% 原书页码：261；段落
这封奏折很简短，两三句说了地方上的问题，就转人了禁教的话%
题。北京传教士得知后，迅速派人了解事件的原委，了解什么原因导%
致福建官员突然要求禁教。传教士探询一番以后，发现并不是什么大%
事。当时福建有两位西班牙新来的传教士，没有遵循男女有别的原则，%
在福安允许男女同时出现在教堂的宗教仪式里面。为什么闽浙总督满%
保要把这件地方上不算大的事件捅到朝廷里去，传教士并没有找到原%
因。传教士只打探到满保先是与雍正沟通过，得到雍正的批复以后，%
才又用公开的朝廷渠道，以地方疆臣的身份上了折子给朝廷。现存的%
满语奏折确实印证了传教士的说法，是雍正指挥满保把事情闹到朝廷%
系统里面去的。\NoteRef{7}{16}%
\par
% 原书页码：261–262；段落（跨页续段）
打探中，传教士还了解到福建地方上对修女问题有意见。当时福%
建有遗弃女婴的习俗，传教士因此收容了不少女婴。收容的女婴长大%
后，传教士把她们培养成修女。当时福建以及浙江很多地方男女比例%
失衡，由于男子娶不到老婆，这些地方还存在卖妻、典妻的习俗。福%
建乡绅看到修女，认为这是天主教蛊惑妇女人教，影响了当地的传宗%
接代。他们并不了解这些修女大部分是由弃婴长大而来。%
\par
% 原书页码：262；段落
而福建地方官场上有禁教的说法起源于福安一个秀才皈依了天主%
教。这让福安知县不安，上报了闽浙总督满保，说传教士在当地兴建%
教堂，蛊惑人心。福安知县知道康熙是容留传教士的，也知道传教士%
手里有印票。但福安知县认为此印票是通融传教士在华居住之用，不%
是允许传教士私建教堂用的。福安知县提到，大清明文规定和尚道士%
不得私建庙宇，而传教士理当不能私建教堂。福安知县巡查，整个福%
安县有多处新修教堂。一个小小县城，就有15座之多（满保后来上%
的折子中提到了这一点）。满保让知县利用当地宗族势力，小心处理，%
限制天主教的发展。\NoteRef{7}{17}%
\par
% 原书页码：262；段落
北京神父在写回欧洲的信件中，坦陈男女区分不严格，共同参加%
宗教活动确实是中国的大忌，很容易惹出麻烦。但他们觉得很冤枉，%
因为新来福建的传教士，北京的神父其实都不了解。而且传教士认为%
这个男女区分就算是问题，也只是一个小地方上的问题，全国其他地%
方都对这个问题处理得很好。对于福建当时的状况，波士顿大学的梅%
欧金教授在他研究福建天主教的专著中有过考证，认为当时禁教的折%
子是北京朝廷政治的需要，并不是福建当地问题的反映。\NoteRef{7}{18}一个最直%
接的表现就是雍正颁布了禁教条约以后，反而福建各地并没有严格执%
行，许多地方教会依然长期存在。这也印证了当时耶稣会神父的说法，%
当时福建并没有什么跟天主教相关的大事发生。%
\par
% 原书页码：263；段落
雍正收到福建的折子后，把折子转送到了礼部，让礼部部议。礼%
部给出的最后结论是赞同折子所请，同意禁止天主教。雍正元年的最%
后一个月（1724年1月），雍正发出谕旨禁教。整个过程从满保上奏，%
到最后禁教，事情发展得很快，总共不到三个月。北京传教士从知道%
折子开始，就在动用他们多年在北京建立的关系网，四处活动。他们%
先是找到钦天监里和满保是世交的满人主管，请他写信给满保。同时，%
广州传教士又多次登门拜见两广总督，请他出面帮传教士说情。两广%
总督开始本不想掺和这件事，但抹不开和传教士来往多年的面子，最%
后还是写信请满保三思。满保对各方求情的人，回复都差不多，大意%
就是事情已经捅到北京，后面事态发展他已无力把控。\NoteRef{7}{19}这样婉言回%
绝了各方，欣然事外了。%
\par
% 原书页码：263；段落
之后，北京传教士把活动的精力主要放在了北京。他们毕竟扎根%
北京几十年，在北京官场也是有根基的。他们很快通过各种渠道打探%
到了各种细节。1723年12月29日，巴多明神父从宫里一个内线那里%
得知，皇帝的意思是保留北京的教堂和传教士，但是禁止地方各省教%
堂和天主教传播。\NoteRef{7}{20}巴多明得知雍正这个处理意见的时候，整个案子%
还没有交到礼部商议。这个信息很重要，这让传教士明白禁教就是皇%
帝的意思，调子已经定好了，礼部部议也就是走个过场而已。知道这%
个消息以后，巴多明又找到他以前在内务府里面一个老熟人，从那里%
得知折子送到礼部的时间。\NoteRef{7}{21}%
\par
% 原书页码：263–264；段落（跨页续段）
尽管传教士知道禁教是雍正的意思，但还是决定要在礼部那边做%
最后的努力。他们先是花重金贿赂了一个礼部的小职员，让他帮忙找%
出当年康熙容教的谕旨，原封不动地抄写出几份来。\NoteRef{7}{22}他们希望康熙%
朝的这份谕旨能派上用场。他们先后通过各种关系，找到礼部中多人。%
其中，康熙的皇十二子、礼部满人尚书以及另外两位官员都同意帮他%
们仔细过问案情。1724年1月3日，礼部部议了满保的折子。会中，%
一位官职低、平时不起眼的官员给出了处理意见，同意满保所请。答%
应帮助传教士的几位，相继在会中发问：为什么整件事都没有参考过%
康熙容教的谕旨？几个官员要求那位下级官员重新思考，下一次部议%
中再给出新的意见。奇怪的是，那位下级官员竟然一字不改，把前一%
天的处理意见再次拿到了第二天的部议上来，而且强调他不会做更改。%
按照传教士的记录，礼部官员从这位官员不同寻常的强势中，察觉到%
了他背后有人，而且很可能就是皇帝。所以这位平时不起眼的官员才%
会在这个跟他私人利益没有关系的问题上公然顶撞部门内的领导。在%
这种情况下，大多数礼部官员都很识相，默默地同意了那位低阶官员%
给出的议案。倒是传教士没有打点过的理学名臣、当时任汉人礼部尚%
书的张伯行，以及另外一名与传教士有交往的礼部官员，拒绝在议案%
上签字。不过二人因为此事受到很大压力，几天后还是同意了议案。\NoteRef{7}{23}%
\par
% 原书页码：264；段落
这个时候，传教士看礼部方面爱莫能助，也确认了禁教就是雍正%
的意思。由于他们已经被赶出了内务府，早已不可能像以前见康熙那%
样，去直接求见雍正。他们就去找了雍正朝唯一可能帮到他们的人，%
那就是康熙的皇十三子。传教士以前和康熙的各位皇子都有交往，知%
道各位王爷的喜好。他们知道皇十三子喜欢书画，而且特别喜欢传教%
士郎世宁。郎世宁在康熙晚年才到北京，在传教士中属于晚辈，出%
面应酬的事一般还轮不到他。这次两位神父特地带着他一起去求见皇%
十三子。%
\par
% 原书页码：265；段落
雍正登位第一年，最忙碌的臣子大概就是皇十三子。传教士第一%
次上十三王爷府，皇十三子正要出发入宫，但还是匆忙中跟他们说了%
几句。这也可见传教士确实在康熙朝积累了过硬的关系。当时想要求%
见皇十三子的达官贵人太多了，很多都排不上号。皇十三子一见到他%
们就说，知道你们是为什么事情来的，告诉他们这事已经转到他和皇%
十六子这里处理了，他还没有具体了解过整件事，让传教士第二天到%
朝廷那边去见他。当时传教士感觉皇十三子的口气和善，似乎看到了%
希望。%
\par
% 原书页码：265；段落
第二天传教士去朝廷那边等了一天，都没有见到皇十三子。那时%
年羹尧正在西北平定罗卜藏丹津叛乱，战事吃紧，皇十三子当天公务%
繁忙，让传教士明天到他府上去商谈。传教士第二天一直等到晚上六%
点才终于见到了皇十三子。皇十三子过问了这件事，已经知道禁教是%
雍正的意思。于是一来就跟他们说事情很棘手，大意是他也帮不上%
忙了。\NoteRef{7}{24}%
\par
% 原书页码：265；段落
第二天传教士们写了一份陈情书交给皇十三子。送陈情书的太监%
出来告诉他们皇十三子和皇十六子看了他们的上书。晚上他们等到了%
皇十三子，皇十三子说他认为陈情书措辞不恰当，整篇文字是在和皇%
帝争辩的架势，这样只会起到反作用。同时说如果他们实在要把这陈%
情书交到皇帝那里，要有什么不好的结果就不要再来找他了。\NoteRef{7}{25}%
\par
% 原书页码：265；段落
听完皇十三子的话，传教士感觉情况严重，也就没有提交他们的%
陈情书。皇十三子同时还告诉他们，跟皇帝上书要以谢恩和祈求的方%
式来写，这样才会有结果。还没有等到传教士再次写陈情书，雍正禁%
教的谕旨就下达了，而且很快就知会了地方。%
\par
% 原书页码：266；段落
传教士面对这样的局面已经无计可施。几个月后，他们又再次去%
找皇十三子，告诉皇十三子各省传教士在禁教后遇到的具体困难，希%
望皇十三子能帮他们疏通。按照雍正禁教的安排，各省传教士要被驱%
逐去澳门。但当时澳门被葡萄牙控制，许多非葡萄牙籍的传教士到那%
里反而更难找到回国的船。所以传教士希望朝廷能允许他们去广州，%
那里各国来的船都有。同时许多地方上的传教士年事已高，外加在当%
地都住了几十年，突然要离开，需要更多的时间来善后。皇十三子听%
完这些具体困难，认为雍正会同意他们的请求，然后教他们应该如何%
上书雍正来提这些问题。皇十三子告诉传教士，禁教之事，木已成舟，%
新上折子就不要再对禁教之事做争辩。同时，皇十三子教传教士折子%
前半部分强调传教士人华百年来奉公守法，然后祈求皇上开恩，宽限%
一些时日。最后皇十三子交代他们把折子送到内阁，让大学士帮忙转%
呈皇上。同时皇十三子告诉他们如果那边不能把折子递上去，再来找%
他，他用他的渠道帮忙送呈皇上。\NoteRef{7}{26}%
\par
% 原书页码：266；段落
上面皇十三子的指点是传教士写回欧洲的信中提到的。由于传教%
士这封折子是通过内阁提交上去的，因而该折子得以保留在宫廷档案%
中。折子全文如下：%
\par
% 原书页码：266；段落
{\bookkai 西洋人戴进贤等谨奏，为吁恩垂鉴事切。臣等自利玛窦航海东来，}%
{\bookkai 历今几二百年。幸荷圣朝优容无外，故士至如归，守法焚修，原非左}%
{\bookkai 道。兹因福建之事，部议波及各省，一概驱往澳门。远臣奉命惟谨，}%
{\bookkai 敢不凛遵。}%
\par
% 原书页码：266–267；段落（跨页续段）
{\bookkai 惟是澳门非洋船常到之地，若得容住广东，或有情愿回国者，尚}%
{\bookkai 可觅便搭船；今俱不容托足，则无路可归。澳门虽住洋商，而各省远}%
{\bookkai 臣，不同一国者甚多，难以倚靠。可怜欲住不能，欲归不得，此诚日}%
{\bookkai 暮途穷之苦也。}%
\par
% 原书页码：267；段落
{\bookkai 近接广东来信，抚臣奉文之后\AuthorNote{〔指朝廷禁教条文〕}，出示行牌，严}%
{\bookkai 加催逼，限六月内驱往澳门，不许迟过七月。因思臣等荷蒙圣恩，留}%
{\bookkai 京备用，则每年家信往来，亦所不免。倘广东无人接应，将来何以资}%
{\bookkai 生。我皇上仁恩溥博，薄海内外咸荷覆帱\AuthorNote{〔受到恩惠意〕}，似此老迈孤}%
{\bookkai 踪，栖身无地，不得不冒渎严威，惟望圣恩宽厚，俯赐矜全，行令广东，}%
{\bookkai 免其驱逐。嗣后各省送往之西洋人，愿赴澳门者，听往澳门；愿住广}%
{\bookkai 东者，容住广东。如此则臣等感激涕零，受恩靡尽矣。再各省现有衰}%
{\bookkai 老病废难行之人，可否暂容？此又出自皇上格外隆恩，非臣等所敢擅}%
{\bookkai 请也。臣等不胜呼号待命之至。谨缮折具奏，伏乞皇上睿鉴，特赐俞}%
{\bookkai 允施行。\NoteRef{7}{27}}%
\par
% 原书页码：267；段落
很明显，折子就是皇十三子提的思路：开篇说传教士一直遵纪守法，接%
着讲困难，最后祈求皇上开恩。传教士虽然很多懂中文，但这篇折子，%
遣词造句很可能是找过汉人帮忙润饰修改而成，比他们其他的中文折%
子高出一个水平。最后“臣等不胜呼号待命之至”“伏乞皇上睿鉴”，把%
他们趴在地上求皇上开恩的即视感都烘托出来了，深得文臣奏议精髓，%
神情并茂地把皇十三子交代的“祈求皇上”的意思写了出来。%
\par
% 原书页码：267–268；段落（跨页续段）
这个折子大概也真是感动了雍正。根据传教士记载，雍正接到奏%
折后立马就给了批复。当时去送折子的三位传教士还在跟大臣寒暄，%
尚未离开。官员接到送出来的雍正回复以后，看了一眼，马上就把批%
复递给了传教士。巴多明神父看完回复后，马上用满文说了一句官场%
逢迎的套话。巴多明到底说了什么话，传教士的记录没有写明。他们%
只记录了传批复的仆人，一听到巴多明的满文回复，马上就跑回去禀%
报雍正。而雍正听完后，立马就召见了三位传教士。这次召见对传教%
士来说意义非凡。自从康熙去世，他们被赶出内务府以后，就没有机%
会见到皇上。这次会见，是自福建上禁教折子以来，快一年时间里第%
一次有机会见到皇上。这跟他们在康熙年间随时能见到皇帝完全是两%
重天。%
\par
% 原书页码：268；段落
传教士记载，雍正大约接见了他们十五分钟，雍正说什么，传教%
士都点头称是，没有任何争辩，只是一再谢恩。感谢雍正同意让传教%
士去广州。雍正说完以后，也对传教士的反应表示满意，还赏了每人%
一件小礼物。\NoteRef{7}{28}%
% 原书页码：268；标题
\subsection{雍正为什么要禁教?}
\par
% 原书页码：268–269；段落（跨页续段）
雍正为什么要禁教？这是天主教历史中一个绕不开的话题。雍正%
登位后第一年就禁教，而这时全国上下并没有什么跟天主教相关的大%
事发生。福安县反映的案子，指控教堂内男女混杂，当时的大臣和后%
来的学者都知道那不过是禁教的托辞罢了。每年无论寺庙还是道观，%
出现这种小范围的问题可以说数不胜数，地方官连朝廷都不会上报，%
就把事情处理了。既然当年没有什么大事，后世学者自然把目光放远，%
把禁教联系到之前发生的大事上去。这也就是为什么提到雍正禁教，%
原因都联系到了康熙时期的礼仪之争和康熙会见教皇特使这样的“大%
事”，甚至还有的把禁教与20世纪以来流行的中西文明冲突理论联系%
到一起。以前的说法，总结起来就是雍正禁教是长时间积累的结果，%
康熙时就有了苗头，雍正元年禁教只是表明雍正做了康熙准备做还没%
有做的事。历史学家宏观看待问题，把事件放到历史长河中去思考，%
这种思考方法本身并没有什么问题，但问题是这样的宏观思考，忽视%
了雍正这个个性鲜明的皇帝自身的作用。稍微了解清史的都知道，雍%
正是个有主见的皇帝，他做决策有他自己的决断思路和逻辑。在禁教%
问题上，雍正才是真正的主角。%
\par
% 原书页码：269–270；段落（跨页续段）
清初皇帝手中的权力跟明朝皇帝完全不一样。康雍乾三个皇帝不%
仅手上都有实权，而且他们每个人也都追求对国家的管控。黄仁宇的%
《万历十五年》很生动地讲述了明朝皇帝被朝廷制度限制，不能放开%
手脚切实地管理国家，以至于万历皇帝二十多年都不上朝，让朝廷各%
部自行运作管理。尽管大清继承了明制，但是大清皇帝有人关前就形%
成的满人权力运作网络。皇帝如果遇到在朝廷那边不能按心意解决的%
问题，可以直接绕过朝廷处理，最直观的就是康熙雍正都绕开朝廷，%
有自己的秘密通信系统。皇帝可以直接通过主奴关系，秘密地和地方%
上的满人联系，绕开朝廷管道来办他们想办的事。比如前章讲到的康%
熙两次接待教皇使团，全部过程没有经过任何朝廷部门，从广东一路%
到北京，康熙自己的内务府就一手处理了。朝廷的史官，无论是《康%
熙起居注》还是后来的《康熙实录》，一条康熙接待两个教皇使团的%
记载都没有，就是这个原因。有些事需要朝廷部门配合的，皇帝秘密%
联系了他在地方的官员后，再让官员把事情递交到朝廷里处理，这样%
里应外合，事情很容易就按照皇帝的想法来办理。前面讲到的闽浙总%
督满保的折子就是这个情况。满保本来只是私下向雍正汇报了一些地%
方上天主教的情况。雍正表扬了满保“卿所言甚是”以后，御批最后%
一句让满保“如此缮本具奏”，\NoteRef{7}{29}就是让他换成正式官文样式重新把事%
情公开上奏到朝廷里面去。后来朝廷礼部的部议其实就是走个过场。%
以前历史中也注意到雍正逐渐把权力集中到他自己的手里，称之为加%
强“中央集权”，其实更准确的说法应该是“个人集权”。明代也是%
中央集权，但那是把权力集中在朝廷中央，而不是皇帝个人。清朝康%
雍二帝利用满人的渠道把权力集中到了个人，所以在看待康熙雍正朝%
的事件时，得把他们个人意志放在首位。%
\par
% 原书页码：270；段落
历史中的细节也印证了禁教是雍正自已要禁教，根本没有什么历%
史趋势以及大的历史事件迫使雍正禁教。在知道福建递上来的折子后，%
康熙的十二子、十三子、十六子在传教士找到他们后的第一反应都是%
他们可以帮忙。但他们最后没有帮上忙，就是因为知道那是皇帝个人%
的意思。但他们的第一反应都真切地说明当时并没有什么大事，几个%
皇子都没有感受到有什么原因需要禁教。这几个皇子都是雍正的人，%
反应尚且如此。其他跟传教士交好的，如皇三子、九子、十子，这时%
都已被流放。当时跟雍正关系最近，也是雍正朝权力最大的皇十三子，%
就当着传教士和皇十六子的面说过，他一看到满保的奏折就觉得不妥，%
认为不能拿一个小地方事件，做成一个全国范围内的政策。而且礼部%
中的满人和汉人尚书在第一次会议中，同时都反对这个折子所请。\NoteRef{7}{30}%
虽然他们二人最后都妥协了，但他们最初的反应也说明当时社会没有%
什么大家公认的“潮流”“趋势”在支持禁教。这些历史参与者的第%
一反应是值得注意的。%
% 原书页码：271；标题
\subsection{雍正对禁教的解释}
\par
% 原书页码：271；段落
上面提到北京传教士求雍正开恩，让各地传教士可以去广州的折%
子。对于这个折子，雍正给了回复。这个回复应该是雍正第一次正式%
向传教士解释他为何禁教。他的朱批是这样写的：%
\par
% 原书页码：271；段落
{\bookkai 朕自即位以来，诸政悉遵圣祖皇帝宪章旧典，与天下兴利除弊。}%
{\bookkai 今令尔等往住澳门一事，皆因福建省住居西洋人在地方生事惑众。朕}%
{\bookkai 因封疆大臣之请、庭议之奏施行政者，公事也。朕岂可以私恩惠尔等，}%
{\bookkai 以废国家之舆论乎？今尔等既哀恳乞求，朕亦只可谕广东督抚暂不催}%
{\bookkai 逼，令地方大吏确议再定。\NoteRef{7}{31}}%
\par
% 原书页码：271；段落
由于传教士的折子不是以私人的密折形式递交给雍正，而是通过内%
阁传递，这样一来，雍正的回复就是公文性质。他知道他的回复内%
阁里的人都会看，所以这个回复有很明显的官腔。回复的实际意思%
是同意晓谕广东不要催逼传教士，但前面一大部分是在为其禁教做%
解释。雍正一开头就声称他继位以来，一直遵循康熙的政策。而禁%
教是“封疆大臣之请”，结果是大臣“廷议”的结果。这样一来，雍%
正把自己在里面的角色撇得干干净净，而他只是在根据朝臣意思办%
公事”罢了。%
\par
% 原书页码：271–272；段落（跨页续段）
雍正在一封给两广总督孔毓珣的御旨中还解释过他为何禁教。孔%
毓珣于1724年4月出任两广总督。他是山东曲阜人，据称是孔子的%
后代，当然也是汉人。他接替的是杨琳。杨琳是八旗里面的汉军旗人。%
在旗人内部，汉军旗人一定程度上是被当满人看待的。杨琳给康熙上%
奏折都是以奴才自称。他在广东任上七年，帮助康熙处理了很多西洋%
事务，跟各路洋人传教士都有私交。杨琳被调离两广总督职位后，雍%
正没有给他新的任用。孔毓珣上任后向雍正汇报了他了解到的西洋人%
以及传教士的情况。他上任之时，禁教谕旨已经颁布了三个月。他说%
他调查广东各地后，并未见传教士有违法乱纪之事。他也不知道到底%
雍正禁教的目的是什么，只能附和礼部的处理意见。雍正收到孔的上%
奏后，给了他回复，回复是这样写的：%
\par
% 原书页码：272；段落
{\bookkai 朕不甚恶西洋之教，但与中国无甚益处，不过从众议耳。你酌量，}%
{\bookkai 如果无害，外国人一切从宽好。恐你不达朕意，过严，则又不是矣。\NoteRef{7}{32}}%
\par
% 原书页码：272；段落
这封回复中，雍正让孔毓珣从宽对待传教士之前，简单解释了他禁教%
的原因，他是“从众议”。这里雍正又把禁教归结为朝廷廷议结果，%
说明这不是他的个人意愿。%
\par
% 原书页码：272；段落
为什么雍正在两次解释中都把禁教原因推到朝廷众议上面？原因%
可能很简单，因为他根本就没有一个令人信服，特别是让这些朝廷高%
官信服的理由。如果天主教或者传教士当时真有什么问题，他完全可%
以直接坦白地告诉他的朝臣，不用把禁教的决定推到“众议”上。还%
有一点，雍正反复说明禁教不是他的个人意思，其实正说明当时的官%
场就是这样认为的，所以雍正才感觉有必要反复说明。%
% 原书页码：273；标题
\subsection{当局者怎么说?}
\par
% 原书页码：273；段落
关于禁教，雍正有他的说法，其他人也有他们的理解。禁教不久，%
当时在北京的耶稣会神父就详细记录了他们了解到的全部过程，其中%
就有他们自己对禁教的分析。他们认为传教士被北京权贵们疏远，到%
最后的禁教，都和雍正的个人喜好有关，他们是这样说的：%
\par
% 原书页码：273；段落
{\bookkai 时过境迁，今上几乎不再使用欧洲人，对科学及外国珍奇之物似}%
{\bookkai 乎也鲜有兴趣。皇帝这种态度使原先一些朋友疏远了我们，他们中一}%
{\bookkai 些人已不能再为我们出力，另一些则不敢与外国人来往。\NoteRef{7}{33}}{\bookkai 33}%
\par
% 原书页码：273；段落
雍正确实和康熙在个人兴趣爱好上有天壤之别。尽管这些喜好是个人%
的偏向，但这个偏向决定了他们的待人接物的态度。朝臣权贵们都是%
明眼人，特别是雍正一登基就把传教士赶出了内务府，这已经释放出%
了明确信号：雍正对这些传教士不感兴趣。在这种情况下，传教士的%
处境自然急转直下。在传教士理解中，福建发上来的禁教奏折就是整%
个官场环境的一个反映。%
\par
% 原书页码：273；段落
皇十三子和传教士原本算不上深交。在康熙年间，传教士的书信%
中基本没有提到过皇十三子。尽管如此，皇十三子作为雍正朝初年最%
有实权的王爷，还是尽其所能帮助了传教士。根据传教士所记录，皇%
十三子和其他几位皇子曾多次接见他们。%
\par
% 原书页码：273–274；段落（跨页续段）
毕竟这些传教士在康熙内务府里几十年，跟这些皇子们家里家外%
都是几十年的老相识。传教士性格都很好，皇子们自已或者家里人生%
病，还会经常找他们看病。在这些皇子们眼中，传教士都是一些老好%
人。所以当传教士找到皇十三子，皇十三子没有敷衍他们，也到雍正%
那里帮他们疏通过。后来皇十三子还不厌其烦地指导传教士如何给雍%
正上折子。传教士写好后，他还帮他们修改。特别需要注意，雍正初%
年皇十三子是最忙的王爷，传教士在他们的记录中也多次提到他很忙。%
就是在这样的背景下，皇十三子还在尽力帮传教士，可以说是仁至义%
尽了。从传教士的文字中看，他们也确实很感谢皇十三子的帮助。在%
一次闲谈中，皇十三子还帮传教士做了分析，也特别告诉传教士那是%
他个人的理解：%
\par
% 原书页码：274；段落
{\bookkai 你们不要弄错，这是我的想法。你们知道今上在只是四亲王时是}%
{\bookkai 很喜欢和尚道士的，不过当初他尚未登基。你们要听明白，我的先父}%
{\bookkai 皇对你们十分喜欢，恩宠有加，但你们不是不知道，他对你们的庇护}%
{\bookkai 与恩赐常引起文人们窃窃私议。\NoteRef{7}{34}}%
\par
% 原书页码：274；段落
这段皇十三子私下的话，是要让传教士知道，你们当初得宠是康熙个人%
喜欢你们。同样，雍正喜欢和尚道士，所以你们就失宠了。皇十三子通过%
这一最直白的对比来劝慰传教士，让他们想开点。言外之意，这些个人%
喜好，是没有办法去找关系疏通改变的。这就是命，让他们认了，不要再%
折腾找人了。从后世的历史角度看，这段话其实涉及一个很重要的历史%
问题，到底这些有实权的皇帝的个人喜好对历史有多大影响？皇十三子%
提到的雍正喜欢“和尚道士”，这个“喜欢”在历史中有多大分量？%
% 原书页码：275；标题
\subsection{雍正与佛教}
\par
% 原书页码：275；段落
雍正信佛，后世一直都知道，不是什么新闻。信佛的皇帝在历史%
上有很多，但真正让雍正在历史上独树一帜的是他“信”的程度。简%
单来说，雍正不是大众认知中一般烧香拜佛的佛教信徒，而是一个闭%
关打坐的修行者，而且达到了佛教内部都认可的一流境界。可以这样%
说，普通佛教信徒和雍正的差别，就像一个偶尔打打乒乓的人和国家%
队一线球员的差别。说雍正是佛教信徒，这没有错，但是这个说法使%
后世大大低估了雍正的佛教水平，同时也淡化了佛教对雍正的影响。%
\par
% 原书页码：275；段落
说到雍正信佛，欧美学者认为清初皇帝都宣称崇信佛教，这是他%
们的政治需要，因为佛教信仰能帮他们更有效地统治广大的蒙古和西%
藏地区。因而信佛是一种政治手段。\NoteRef{7}{35}民国清史大师孟森尽管深知雍%
正信佛修行，但他基于雍正上台后残酷镇压异已和兄弟的历史，不相%
信雍正是真正的佛教大师，因而解释雍正是通过佛教修行来掩盖他%
夺位的企图。\NoteRef{7}{36}中外史家的这些解释，都混淆了一个时间顺序，那就%
是雍正先是佛教修行者，然后才是皇帝。雍正在他20岁出头就与佛%
教结缘。雍正还是皇子的时候，距离他的藩邸（现在的雍和宫）不到%
三百米就是当时著名的柏林寺。这对雍正亲近佛法可谓是得天独厚的%
条件。后来，曾做过柏林寺方丈的独超禅师圆寂后，其弟子来找雍正%
为其写塔铭。其中就写到二人是康熙四十一年开始一起谈经说法的。\NoteRef{7}{37}%
\par
% 原书页码：275–276；段落（跨页续段）
康熙四十一年（1702年），雍正不过二十五岁，有缘拜会当时的%
一代名僧，还能与其“谈法甚契”，这说明雍正学佛已经有些时日了。%
而且雍正称独超为“师”，这些前后因缘都说明雍正跟独超禅师学过%
佛，不是一般的应酬关系。%
\par
% 原书页码：276；段落
雍正自号破尘居士、圆明居士。他不是历史中常见的那种只“谈”%
佛法的文人，而是一个参禅打坐的实修者。他自己记录了他参禅中破%
参的经历，写道：%
\par
% 原书页码：276；段落
{\bookkai 壬辰春正月，延僧坐七、二十、二十一，随喜同坐两日，共五枝}%
{\bookkai 香，即洞达本来。方知惟此一事实之理，然自知未造究竟。……叩问}%
{\bookkai 章嘉，乃曰：“若王所见，如针破纸窗，从隙窥天。虽云见天，然天体}%
{\bookkai 广大，针隙中之见，可谓遍见乎？佛法无边，当勉进步。”\NoteRef{7}{38}}%
\par
% 原书页码：276；段落
壬辰年是1712年，这年春，康熙的太子还在位。这说明雍正修佛先%
于后来的储位之争。雍正这段记录有点今天所谓“炫技”的意思，一%
起同修的和尚要打坐七天或者二十、二十一天，而他只打坐了五枝香%
的时间就“洞达本来”。雍正记录提到，他的参悟还得到了章嘉活佛%
的认可，确认他是入门了，让他继续。雍正还记录了他在禅修中参破%
三关的体验：%
\par
% 原书页码：276；段落
{\bookkai 二月中，复结制于集云堂，着力参求。十四日晚，经行次，出得}%
{\bookkai 一身透汗，桶底当下脱落，始知实有重关之理。……恰至明年癸巳之}%
{\bookkai 正月二十一日，复堂中静坐，无意中忽蹋末后一关，方达三身四智合}%
{\bookkai 一之理，物我一如本空之道。庆快平生。诣章嘉所礼谢，国师望见，}%
{\bookkai 即曰：“王得大自在矣。”\NoteRef{7}{39}}%
\par
% 原书页码：277；段落
禅修中有破三关的说法，是禅修的三个阶段。雍正上面三段记录就是%
分别写出了他破三关的主观体验以及时间。而且他的破关经历还得到%
了章嘉活佛的确认。这位章嘉是章嘉活佛第一代（1714年圆寂），是%
当时著名的高僧。\NoteRef{7}{40}在佛教传统中，一个人禅修中的证悟，是否修行%
得法，以及是否误人魔道（歧途），这些都是由在世高僧来确定。所%
以雍正的记录中，用章嘉活佛的话来确认他的修行，符合佛教传统。%
从历史角度看，雍正这些记录的可信度很高，因为这些记录出自他编%
辑的《御选语录》。这本书是历代高僧语录的集结。雍正的本意是希%
望用他编出的语录来帮助后来的人修行。因而雍正很明确，他编的这%
本书的读者是佛教中有一定修为的修行者。雍正需要向这些人说明白，%
凭什么他有资格来编辑这些高僧语录。他记录下他自己破三关的经历，%
其实就是在向后来的修行者展示他的资格。%
\par
% 原书页码：277；段落
到底雍正的佛教修行水平如何，没有修行过的人当然无从判断。%
他所写的主观经历，没有经历过的人也无法辨别。要评鉴他的水平，%
现在流行的同行评议原则大概可以作为一个比较公允的手段。首先我%
们可以看雍正自己对自己水平的一个评估，然后再看佛教中公认的高%
僧对他的评价。雍正自己曾经评论过禅修中的三关，他说：%
\par
% 原书页码：277；段落
{\bookkai 三关妙旨，累进而上，言思俱绝，历来古德宗匠，多止提倡初步。}%
{\bookkai 至于重关，不过密隐于语句之中，令至者自会，未尝轻以一字一句，系}%
{\bookkai 缀学人，累他堕入情推识解也。即观向来宗师透重关者，亦不多见，近}%
{\bookkai 年以来，几至埋没不续矣。朕因加意提撕，而诸人感朕开示之恩，努力}%
{\bookkai 精进，于是因真参而得实悟透重关者，颇有其人，此实从来罕遇者。\NoteRef{7}{41}}%
\par
% 原书页码：278；段落
按雍正解释，历史上的高僧，一般都只在文字中记录下第一关（通常%
称为“本参”）的修行。而对第二关的重关没有多加解释，留给破这%
一关的人自行领悟。这就是雍正所谓的“至者自会”；同时他解释高%
僧们这样做的原因是怕立下文字以后，误导了后人。因为有了文字，%
后人就会通过文字，用逻辑理解的方法来思考重关，就是雍正所谓的%
“累他堕入情推识解”。这样的话，反而误导了后人的修行。这和佛教%
中著名的“以手指月”的故事暗合。佛教中认为文字就像手指，本来%
是指着月亮的，但不知道的人，就会围着手指琢磨，忘记了月亮才是%
手指指向的意义。%
\par
% 原书页码：278；段落
对于“重关”，雍正的评价是：“即观透重关者，亦不多见，近年%
以来，几至埋没不续矣。”按雍正的说法，破重关在历史上已经很少%
见了。而雍正的水平是不仅破了重关，还破了第三关，同时他还指%
导一些人破了重关。从雍正的这段评价三关的文字可以看出，雍正%
是把自已放在中国佛教修行的千年历史中做的评价。他的自我评价%
是很高的。%
\par
% 原书页码：278；段落
现当代史中公认的第一高僧大概非虚云和尚莫属。虚云和尚在%
1953年上海玉佛寺的一场法会上，讲到了雍正。虚云说：%
\par
% 原书页码：278；段落
{\bookkai 雍正帝在皇宫里也时常打七。他对禅宗是最尊重的，同时他的}%
{\bookkai 禅定也是非常得好。在他手里悟道的有十余人，扬州高旻寺的天慧}%
{\bookkai 彻祖，也是在他会下悟道的。禅门下的一切规矩法则，皆由他大整}%
{\bookkai 一番，由是宗风大振，故人材也出了很多。\NoteRef{7}{42}}%
\par
% 原书页码：279；段落
1953年虚云已至晚年。虚云提到雍正，因为他当时正在指导僧人坐禅%
七，禅修法门和雍正当年所用一样。因此他在禅七的第二天讲到了雍%
正。虚云不仅肯定了雍正的“禅定也是非常得好”，而且还提到雍正%
指导了多人悟道。虚云提到的天慧彻祖，是扬州高旻寺第一代祖师，%
他是在雍正的宫里禅堂中，在雍正指导下悟道的。佛教中自悟已经不%
易，还能指导多人悟道，这在佛教史中也是不多见的。虚云几句话已%
经把雍正放到一代宗师的位置上了。%
\par
% 原书页码：279；段落
另外，著名的印光法师（1862—1940）在雍正编的佛教书籍重%
新出版时，为雍正作注。其中印光法师确认雍正是涅槃觉悟，对雍%
正的书评价道：“言言见谛，语语归宗。如走盘珠，似摩尼宝。凡具%
眼者，无不佩服。”\NoteRef{7}{43}如果说这些大师的评价还有一定主观性的话，那%
么客观上看，雍正编的佛教书籍后来被中日两国佛教界都收人了佛%
藏中。20世纪，海峡两岸的佛教出版社都在纷纷重新出版印刷雍正%
的佛教书籍。%
% 原书页码：279；标题
\subsection{作为佛教居士的雍正}
\par
% 原书页码：279–280；段落（跨页续段）
雍正从年轻时开始佛教修行，到1712、1713年参禅破关，可以说%
他从十几岁到三十几岁的生活重点都在佛教修行上。他破关之时离他%
登基都还有十年时间。也就是说他在当皇帝之前就是当时佛教界的风%
云人物了。他破关后不仅指点其他和尚修行，还参与到了禅宗内部的%
一些理论争论中。佛教有不立文字的传统，认为文字是一种束缚。所%
以关键问题的解释权都在悟道之人那里。雍正的禅修得到章嘉活佛的%
认可以后，雍正在佛教内部也就有了话语权，是权威。加上后来还有%
修行者在他点拨下破关，这就进一步加强了他在佛教界的权威地位。%
从历史角度看，雍正的权威感首先来自于佛教界，是在精神领域，而%
后来做皇帝带来的权威是在世俗层面。%
\par
% 原书页码：280；段落
雍正在佛教界，用现在的话来说是信仰精神领域的权威。这是一%
个被长期忽略的问题。顺治信佛，后来还曾剃度，但顺治在佛教内部%
看来就是一个普通的信徒，不是精神上的权威。康熙尽管征战四方，%
被蒙古、西藏很多地方视为神，但他自己清楚他在任何宗教里都不是%
精神上的权威。他和教皇谈判就是明证，他清楚教皇才是教内的权威，%
他只是作为一个世俗皇帝在和教皇谈判。但雍正不一样，雍正是佛教%
里被认可的觉悟之人，而且是指导当时高僧修行的人，因而他自信有%
资格评论世间宗教。除了本章开篇引用雍正在得知葡萄牙使团入华以%
后，评论天主教的“天主”概念以外，他在接见葡萄牙使团之时，当%
着传教士的面，还评论道：%
\par
% 原书页码：280；段落
{\bookkai \AuthorNote{〔根据传教士记载翻译〕}你们批评其他宗教教派。难道是一些宗教}%
{\bookkai 好，一些不好？事实并非如此，所有宗教到最后都是殊途同归的，都}%
{\bookkai 在某方面是好的。但从众多宗教论说来看，文人们\AuthorNote{〔指儒家〕}还有你}%
{\bookkai 们这些天主教传教士以及其他一些人，很少有真正把自己宗教教义搞}%
{\bookkai 明白的。你们得先把自己的本心弄明白了，才能真正理解你们自己的}%
{\bookkai 教义。\NoteRef{7}{44}}%
\par
% 原书页码：280–281；段落（跨页续段）
雍正当时已经是皇帝了，他敢于直接评论天下一切宗教，还能说出“宗%
教到最后都是殊途同归”这样的话，这就是雍正作为佛教“觉悟”之%
人在对众生指点迷津的姿态。他还直接告诉面前的传教士，他们自己%
都没有搞明白自己的教义。从传教士记录的文字来看，他们并不认可%
雍正的话，他们记录的本意其实是想体现出雍正的自大和荒唐。但传%
教士没有明白的是，雍正的话其实是一个佛教觉悟者的正常评论方式。%
他的这种评价方式在禅宗公案、佛教故事中，可以说是层出不穷。因%
而，了解了雍正在当时佛教界的地位，就会理解上面评论反映出来的%
不是自大，而更多的是佛教已经融会到雍正的思想灵魂深处。%
% 原书页码：281；标题
\subsection{皇位争斗与禁教的关联}
\par
% 原书页码：281；段落
雍正和传教士之间并没有因为信仰问题有过矛盾或者争论。在康%
熙的引荐下，传教士和他的皇子们都有往来，雍正也不例外。但雍正%
和传教士处于道不同不相为谋的状态，双方没有深交。笼统而言，天%
主教是一神教，本质上传教士不能兼容其他的宗教。传教士能和儒家%
文人保持良好关系，因为儒家实际上不是一个宗教。传教士人华以后%
没有与和尚融洽相处的例子。这一点和中国传统中和尚道士能够互%
相交流不一样。而且在传教士的文字中，还经常看到对和尚的奚落和%
嘲讽。%
\par
% 原书页码：281–282；段落（跨页续段）
天主教不能容纳接受其他宗教这一基本点，注定了传教士不可能%
和皇子时的雍正建立亲密关系。毕竟雍正的府上常年香烟缭绕，和尚%
喇嘛不断。也是因为这个原因，雍和宫后来成了寺庙，现在香火依然%
很旺。而且康熙在世之时，并没有特别迹象显示学佛的皇四子可能继%
位，这样传教士也就没有特殊的理由要进一步去接触皇四子。要是雍%
正不继位的话，他就是康熙十多个皇子中的一位，如此而已。总而言%
之，在雍正登位以前，他研习佛法这一点，提前注定了传教士和雍正%
之间的距离，二者不可能有比较亲近的关系。在康熙晚年长达十多年%
的储位斗争中，传教士没有在雍正这边下过功夫。雍正登位以后，也%
就不把传教士当自已人。重组内务府的时候，把传教士赶出宫去也就%
是自然而然的事了。%
\par
% 原书页码：282；段落
有一个历史细节值得注意，就是跟雍正唯一有私人关系的传教士%
德里格在康雍两朝的不同际遇。德里格在康熙宫廷里做音乐。教皇第%
二位使节在北京时，德里格私下告知使节康熙可以允准教皇的禁约，%
这让康熙震怒，以“德里格乃无知光棍之类小人”开头整整写了一页%
纸朱批来批德里格的种种不是。\NoteRef{7}{45}康熙后来把德里格软禁在北京的一%
个教堂里面。德里格不是耶稣会的神父，他是罗马传信部派往中国的。%
尽管他和耶稣会神父都在康熙手下做事，但双方关系很差，经常敌对。%
公正地说，德里格在音乐方面是有大才的，他在紫禁城内创作的弦乐%
协奏曲，今天世界各国音乐家都还在演奏，曲风恬淡清新，兼具中西%
特色。康熙爱才，尽管私下对德里格有不满，还是一直把他留在宫中%
组建西式乐团。俄罗斯彼得大帝使团人华，听完康熙的乐团演奏后大%
为吃惊，没想到在紫禁城内还有一支如此高水准的西式乐团。%
\par
% 原书页码：282–283；段落（跨页续段）
对德里格个人来说，他在康熙晚年最大的收获是做过雍正的音乐%
老师，而且一直和雍正保有私交。因为有这个私交，雍正一上台就把%
德里格从软禁中释放了，而且还任用他来负责接待雍正朝所有来华西%
洋使团。也是在雍正上台这一年，德里格终于实现了他多年的夙愿，%
在北京建立一座属于自已的教堂，这样他终于能从耶稣会神父的教堂%
里独立出来。他建立的教堂在北京称为西堂，在现在西直门内大街上。%
德里格的例子就最直观地说明了和皇帝登位前的私交对传教士来说是%
何等地重要。%
\par
% 原书页码：283；段落
雍正信佛决定了他如何看待传教士及他们的用处。雍正和康熙不%
一样，他对传教士讲的数学、天文等欧洲学术都没有兴趣，因而在他%
眼里，传教士并没有什么特别之处。这一点是他和康熙在看待传教士%
问题上的根本区别。%
\par
% 原书页码：283；段落
应该说雍正刚上台的时候，还是有意以一个皇帝的姿态，大度地%
容纳天主教的。传教士穆经远代表葡萄牙国王给雍正送了很多礼物恭%
贺他登基，雍正全部收下了，这在当时来说是一个很好的开端。但是，%
在康熙内务府里混迹了几十年的传教士，不可避免地被牵涉到雍正早%
年的权力斗争中。1723年4月，雍正登基不到五个月，发现穆经远和%
皇九子关系很深，就把他和皇九子一起发配西宁。同时他警告在北京%
的传教士不要多管闲事。\NoteRef{7}{46}%
\par
% 原书页码：283–284；段落（跨页续段）
北京传教士似乎并没有把雍正的警告放在心上。当然也有可能%
是传教士和雍正对不要“管闲事”的理解不一样。雍正元年，北京%
传教士继续和雍正要打击的苏努一家走得很近。苏努是努尔哈赤四%
世孙，与雍正同辈，累世军功，有贝勒封号。在康熙朝，苏努一家%
一直支持皇八子登基。苏努本人没有信教，但是其有四个儿子以及%
数十名家丁从康熙五十一年（1712年）后陆续受洗，都是当时有名%
的天主教徒。因此传教士一直和苏努家有来往。这种来往本来并不%
是什么大事。但是雍正上台后，在关于苏努家的奏报中经常发现传%
教士的身影，这就成了大事。\NoteRef{7}{47}在发配皇九子到西宁时，除了穆经远%
以外，还发配了苏努的两个儿子。雍正一直派人监视皇九子等人的%
言行。但让监视人都感到意外的是，穆经远竟然在1723冬（当年的%
圣诞节）公然给苏努两个儿子受洗，让他们正式成为了教徒。这些%
都被眼线汇报给了雍正。后来，让雍正更气愤的是，他发现皇九子%
为了躲避监视，竟然把西洋文字做成密码在和北京联系。而这西洋%
文字无疑是穆经远教的。用西洋密码联系一事在多个清官方历史中%
都有记载。最近有学者在清宫档案中发现了这些用西洋密码写成的%
私信，还做了总结。\NoteRef{7}{48}%
\par
% 原书页码：284；段落
同时，在北京，传教士私下还和苏努家有来往。这些来往雍正和%
他内部圈子的人都知道。当时雍正圈子中的隆科多就看在眼里，很为%
传教士着急。隆科多是佟国维的儿子，是佟家在雍正朝早期的核心人%
物。前面已经讲过，传教士和佟家有几代的关系。所以隆科多在宫里%
碰到传教士巴多明的时候，冒着风险，告诉他皇帝知道他们和苏努家%
的关系，让他们传教士赶快远离苏努一家。\NoteRef{7}{49}不过，传教士没有认识%
到问题的严重性，在他们看来，他们和苏努家的来往跟政治完全没有%
关系。而且作为传教士，他们很难抛弃追随他们信仰的信徒。%
\par
% 原书页码：284–285；段落（跨页续段）
雍正初年的朝廷政治远比传教士想象中要复杂。巴多明还记录了%
一个细节。隆科多在跟他说话的时候，突然看到迎面另一个大臣向他%
们走过来，便马上停止了谈话。这个细节就是当时朝廷政治的直观写%
照。隆科多是雍正朝头几年最重要的实权派。即便以他这样的地位，%
跟传教士说话的时候都要处处提防，小心谨慎。当时，雍正正在清除%
异己，许多朝臣和皇亲国戚都因为派系问题受到牵连，即使位高权重%
如隆科多都要处处小心。但传教士似乎没有察觉到当时危机四伏，还%
在继续和苏努一家来往。%
\par
% 原书页码：285；段落
禁教和皇位之争有关联，这并不是后见之明的分析，当时跟雍正%
关系最好的传教士德里格，在他1725年（禁教后第二年）发回欧洲%
的一封信中就写道：%
\par
% 原书页码：285；段落
{\bookkai 各省驱逐传教士与礼仪问题完全无关。禁教始于福建，正如我}%
{\bookkai 去年写的那样，皇帝批准了。（有人说他是始作俑者，但我不能确定）。}%
{\bookkai 终究他还是批准了。说他是始作俑者（尽管我不这么认为）源于欧}%
{\bookkai 洲人可能站在他的兄弟那边。也就是说如果他的哪个兄弟要起事，}%
{\bookkai 地方上的欧洲人可以提供很多帮助。\NoteRef{7}{50}}%
\par
% 原书页码：285；段落
德里格尽管也是传教士，但他和当时北京所有的耶稣会传教士都有矛%
盾，他的观点可以看作是独立于耶稣会神父以外的第三方观点。他对%
禁教问题的分析有两点：一是他认为禁教跟康熙时的礼仪问题完全无%
关；二是禁教跟雍正和他兄弟之间的斗争相关，而且这个观点在当时%
就是流行的说法。%
\par
% 原书页码：285–286；段落（跨页续段）
那么到底雍正禁教有多少政治因素在内？这很难知晓。因为就算%
雍正因为政治问题要处理传教士，他也不会直接说出来，让这事留在%
历史记录里。看看年羹尧案，当时抓年羹尧的时候，给出的理由是年%
羹尧的奏折用词不当。后来审讯又给年羹尧列出了一系列罪名，但当%
时人和后来人都明白那些公布出来的罪名不过都是托辞。真正的原因%
一直是个谜，时至今日学者都还在争论。传教士涉及政治这一点，到%
底多大程度影响了雍正的禁教决策，很难有明确的界定。但有影响是%
一定的，否则隆科多也没有必要冒着风险提醒巴多明神父。后来，雍%
正坚持要杀掉穆经远神父也说明他一定是比较深入地卷入了派系斗%
争。雍正知道杀传教士会在历史上留下骂名，但他思虑再三后，还是%
赶在葡萄牙为穆经远神父求情的使团到来前，把他给杀了。穆经远是%
葡萄牙王室派出的神父，他在康熙晚年一直是耶稣会里的领军人物。%
他成为鸦片战争前唯一一个因为政治问题被杀的传教士。但细看现存%
的穆经远审讯记录，他根本没有犯任何真正的罪，有的只是对皇九子%
的忠心。在因为皇九子案被审问时，他坦然面对，没有否认和皇九子%
的关系，他在审讯开头这样说：“我在允禟处行走，又跟随他在西大同，%
前后有七八年了，允禟对我好是人所共知的。”\NoteRef{7}{51}穆经远的罪其实就是%
在错误的时间，错误的地点，跟错了主子。雍正初年，康熙的皇子们%
站错了队的，尚且不能自保，更何况一个传教士。%
% 原书页码：286；标题
\subsection{小结}
\par
% 原书页码：286–287；段落（跨页续段）
其实北京的满人权贵和传教士都知道，当时无论是禁教还是容教%
都不是什么大事，掀不起什么风浪。就是在康熙的最后二十年中，把%
在中国各个教派的传教士加在一起，最多的时候也就140人左右。根%
据人口统计，清初大概有两亿人，所以这个数字放在当时中国，就是%
沧海一粟，无足轻重。\NoteRef{7}{52}相较之下，当时有记录的和尚就有11万，尼%
姑有9000。\NoteRef{7}{53}这几个数字放在一起，140人的传教士基本可以忽略不计。%
按雍正的话说，传教士“人众不过二十”。\NoteRef{7}{54}对这样一个小群体，无论%
是容留还是禁止，可以说完全听凭皇帝的个人意愿。%
\par
% 原书页码：287；段落
历史的各种细节也证明了当时朝廷上下根本没有共同认可的原因%
来禁止天主教。也就是说没有什么后世所谓的历史趋势、潮流来决定%
当时禁教。禁教就是雍正的个人意思。对这一点，当时的王公大臣都%
看得很清楚。禁教风波期间（1723年10月到1724年1月），当时朝%
廷真正的大事是年羹尧出兵西北镇压罗卜藏丹津叛乱（1723年8月到%
1724年春）。叛乱中，有西藏喇嘛直接拿起武器，跳上马背和清军正%
面对抗。对于这些喇嘛叛军，雍正几次命令前线的年羹尧“勿拆毁庐%
舍祠宇，勿扰害庙内番僧”。\NoteRef{7}{55}雍正的这些举动，完全有违清朝军队对%
不投降者格杀勿论的老传统。朝廷上下都不理解，甚至普遍怀疑是雍%
正内务府里的和尚在干扰朝政。后来雍正不得不在朝会上表明，他只%
是在闲暇时与和尚一起谈经说法，他们并没有干预国事。\NoteRef{7}{56}不过，雍%
正想要保护喇嘛和寺庙是他的个人意愿，而清朝军队的士兵却没有雍%
正对佛教的那份情感。他们还是按照清朝军队的老传统，在打胜仗以%
后把所有喇嘛都杀了，同时也捣毁了寺庙。有记载显示，清朝军队一%
次就杀了6000喇嘛。\NoteRef{7}{57}雍正要求禁止天主教和他要保护喇嘛寺庙基本%
是同一时期的事，放在一起就很明显可以看到，雍正的个人倾向在他%
决策中起到的作用。%
\par
% 原书页码：287–288；段落（跨页续段）
对于传教士来说，雍正烧香拜佛，念经打坐，本来和他们是没%
有什么关系的。而康熙驾崩和雍正登位，立马使得雍正的一切都和%
他们的命运关联起来。可惜的是，雍正信佛和他的个人爱好，是传%
教士没有可能改变的部分。他们唯一能祈祷的是雍正能够怜恤他们，%
把他们当成世间各种大宗小教中的一种，不要为难他们。所以传教%
士数次强调他们是遵纪守法的宗教，人华两百年从未给朝廷添过乱，%
话里话外的意思就是希望雍正像对待其他守法宗教一样，一视同仁%
地对待他们。%
\par
% 原书页码：288；段落
但是，传教士忘记了一点。雍正初年，举国上下，没有任何宗教%
团体像传教士那样深深植根在北京的满人权贵圈中。他们不可避免地%
卷人了皇家的家庭恩怨之中。所以，从某个角度说，传教士的命运其%
实是和当时的皇位之争紧密相连的。如果康熙以后继位的不是雍正，%
传教士在华的轨迹完全有可能是另外一番景象。按照传教士的理解，%
这是上帝的安排。按中国的通俗说法，这就是命。%
